\section{Scope, Method, and Qualifications}

\subsection{Reader, question, and claim}

This is a discursive canon-law article for a serious general reader able to follow documentary argument, and secondarily for students of canon law, for clergy and candidates and their families, and for anyone who has to write or speak accurately about the subject. Its governing question: what does the current law of the Catholic Churches oblige of clerics in the matter of chastity, continence, and celibacy; on what authority; with what effects; and how do the Latin and Eastern disciplines differ?

Its claim, argued rather than announced: the three words name three different things---a virtue binding every baptized person by state of life, an abstention from sexual activity, and an unmarried status with three distinct juridical faces---and almost every serious confusion in this field is a collapse of one into another. From that distinction the rest follows: the Latin Church obliges continence and \emph{therefore} celibacy (CIC c.~277 §1), attaches the obligation by a public act before the diaconate (c.~1037), enforces it by a diriment impediment (c.~1087) and a penal norm (c.~1394 §1), and relaxes it by two distinct acts of two distinct authorities (cc.~290--292); the Eastern common law obliges chastity of celibate and married clerics alike (CCEO c.~374), honours both states by name (c.~373), and refers the admission of married men to each Church's own particular law or to special norms of the Apostolic See (c.~758 §3), while agreeing exactly with the Latin Code that no cleric may marry after ordination (c.~804); the two codes' divergence is a Western development rather than an Eastern departure, and CCEO canon 2 requires the Eastern canons to be read from the Eastern legal tradition rather than against a Latin norm; and the Latin Code has left one question genuinely unsettled---whether canon 277 §1's continence binds the married permanent deacon---which no authentic interpretation resolves.

\subsection{Method and source hierarchy}

The article is governed by the repository's articles profile under both its faith-and-theology and its canon-law rules. Concretely, that meant: the controlling text is the official law, never a commentary or an unofficial translation; promulgation and effective dates were verified at the promulgating acts; validity was kept apart from liceity, obligation from permission and recommendation, the ordinary rule from dispensation, derogation, and rescript, and legislation from administrative act; no Latin canon was applied to an Eastern Catholic or the reverse; and each document was cited at its own authority level.

\begin{threestable}{Source class}{Function in the article}{Governing boundary}
CIC 1983 & The Latin discipline, worked canon by canon (sections~2--6, 9, 11). & Holy See Latin web deliveries of Books I, II, IV, VI, with the English canon-range deliveries as identified working aids, all fetched and hashed 2026-07-25; Book VI in the \emph{Pascite gregem Dei} revision. Web deliveries, not \emph{Acta} pages. Delivery defects recorded, not repaired.\\
CCEO 1990 & The Eastern common law (sections~7--9). & Quoted from the promulgation text, \emph{Acta Apostolicae Sedis} 82 (1990), read at thirteen page images rendered from the Holy See's archival volume PDF; one further page read at the text layer and marked. No official English exists; all English is the article's labelled working gloss. The CLSA Latin--English edition was deliberately not used.\\
\emph{Acta} acts (2014, 2015) & The Eastern diaspora norms and the Churches \latin{sui iuris} (sections~7--8). & Born-digital \emph{Acta} fascicles, fetched and hashed 2026-07-25, read at their text layers; page rasters not required because the fascicles are not scans.\\
Conciliar and papal teaching & \emph{Presbyterorum ordinis} 16, \emph{Optatam totius} 10, \emph{Lumen gentium} 29, \emph{Orientalium Ecclesiarum} 5--6, \emph{Sacerdotalis caelibatus}, \emph{Pastores dabo vobis}, \emph{Querida Amazonia} (sections~3, 9, 12). & Holy See Latin and English web deliveries, fetched and hashed 2026-07-25; \emph{Sacerdotalis caelibatus} n.~42 corroborated against \emph{AAS} 59 (1967) 674. Each cited at its own genre: conciliar decree, encyclical, post-synodal exhortation. AAS collation performed only where stated.\\
Dicasterial acts & \emph{Anglicanorum coetibus} complementary norms; the 1998 deacon documents; the collection of authentic interpretations (sections~9, 11). & Holy See deliveries, fetched and hashed 2026-07-25. Executory decrees and administrative acts are marked as such; c.~33 §1 governs their force.\\
Historical acts & Elvira, Nicaea, Siricius, Trullo, Lateran I and II, Trent, the 1917 Code (section~10). & Each at the identified witness named in the References, with the transmission and authority caveats stated at the claim. The 1917 Code was read at a page image of the 1918 Vatican Polyglot printing. Historical practice is never presented as present law.\\
Modern scholarship & The historiographic dispute (sections~10, 13). & Named at literature level from the disputants' standard positions; not consulted at source; not adjudicated.\\
\end{threestable}

\subsection{Included and excluded scope}

Included: the three obligations and their distinction; the Latin discipline of admission, obligation, impediment, penalty, and loss of the clerical state; the Eastern common law and the 2014 diaspora norms; the Churches \latin{sui iuris} and their ranks; the permanent diaconate and the unsettled continence question; the development of the Western discipline at identified acts; the convert-clergy and ordinariate provisions; the magisterial file from the Council to 2020; and the theological arguments on both sides, stated and not adjudicated.

Excluded: the particular law of any diocese, eparchy, or Church \latin{sui iuris}; the proper law of religious institutes and societies of apostolic life beyond the canons that name them; the statutes of the personal ordinariates; the liturgical books, including the rite in which the obligation of canon 1037 is assumed; the New Testament evidence and the patristic corpus except as they appear in the identified legal acts; Orthodox and other non-Catholic law, which governs no Catholic; the empirical literature on clerical sexual abuse; clergy statistics; and every concrete case.

\subsection{Material qualifications}

\begin{enumerate}[leftmargin=1.45em,itemsep=0.25em]
\item Vatican web texts are dated delivery states, not the \emph{Acta Apostolicae Sedis}. Four delivery defects were found and are recorded where they occur: the English Code delivery's ``a\emph{V}ected'' for ``affected''; the Latin Code delivery's ``3 ab impedimento'' at c.~1047 §2; the Latin \emph{Sacerdotalis caelibatus} delivery's \latin{posset conferr} and \latin{eirum}; and the deacon documents' ``OIC''. Others are not excluded.
\item CCEO has no official English version. Every English rendering of an Eastern canon here is the article's own working gloss, made from the \emph{Acta} Latin and marked as a gloss; the Latin governs, and a reader relying on the English for any purpose should return to the Latin.
\item The \emph{Acta} volume PDF for 1990 is a scanned volume with an OCR layer; its quoted canons were read at page images. The 1967 volume's page rasters could not be rendered by the tools available, so \emph{AAS} 59 (1967) 674 is a text-layer reading corroborated against an independent delivery. The 2014 and 2015 fascicles are born-digital and were read at their text layers.
\item The enumeration of the Churches \latin{sui iuris} in section~8 is the article's own reconstruction from a 2010 official information sheet plus two verified 2015 acts. No official online source states the present total. The Serbian exarchate question the 2010 sheet itself raised is left open.
\item The fourth column of the Churches table reports one dicasterial statement of 2014 and no Church's particular law. It is not a statement of any Church's current discipline, and the Eritrean row is expressly outside the 2014 statement's reach.
\item The disputed question of section~9 is reported and not resolved. The evidential weight given to the asymmetry between \emph{Directory} nn.~61 and 62 is the article's own reading of an executory document that cannot in any case change the law.
\item The bounded negative results are bounded: the absence of an authentic interpretation of cc.~277 and 288 is asserted for one dicastery collection as delivered on one date; the silence of \emph{Querida Amazonia} is asserted for one language of one delivery by literal search; the unobtainability of the 1965 Tisserant letter, of the 1997 circular letter, and of a current official enumeration of the Churches is asserted for the routes actually tried.
\item Historical claims about acts not examined at their own witnesses---the 1929--1930 Eastern decrees, the ``Statement \emph{In June},'' the Pastoral Provision, \emph{Ministeria quaedam}---are carried at the level of the documents that cite them, and the article says so at each.
\item Round figures (the number of permanent deacons; the estimate of Ruthenian defections) are reported background, not verified statistics, and no argument here depends on them.
\item The article's own synthesis---the three-term analysis of section~2, the reading of the \latin{ideoque}, the reconstruction of the Churches' count, the ``what would have to change'' analysis of section~13---is labelled as synthesis where it occurs and is attributed to no source.
\end{enumerate}

\subsection{Legal Scope and Currentness}

\textbf{Governing law.} For the Latin Church: the 1983 \emph{Codex Iuris Canonici}, promulgated by John Paul II by the apostolic constitution \emph{Sacrae disciplinae leges} of 25 January 1983 and in force from 27 November 1983, as amended---including Book VI as integrally revised by \emph{Pascite gregem Dei} of 23 May 2021, in force 8 December 2021. For the Eastern Catholic Churches: the 1990 \emph{Codex Canonum Ecclesiarum Orientalium}, promulgated by John Paul II by the apostolic constitution \emph{Sacri canones} of 18 October 1990 and binding from 1 October 1991, as amended.

\textbf{Authoritative language and translations.} Latin governs both codes. The Latin canons were verified at the Holy See's Latin web deliveries (CIC) and at the \emph{Acta Apostolicae Sedis} promulgation text (CCEO). The Holy See's English web delivery of the 1983 Code is an identified translation used as a working aid and is quoted as such; where it diverges from the Latin, the Latin governs and the divergence is reported. The Eastern canons have no official translation and are glossed by the article.

\textbf{Jurisdiction and persons.} Latin conclusions bind only the Latin Church (CIC c.~1); Eastern conclusions bind only the Eastern Catholic Churches (CCEO c.~1). Nothing here may be transferred across that line. The persons addressed by the norms discussed are candidates for and recipients of sacred orders, their spouses where they have them, and the authorities named in the canons.

\textbf{Universal, particular, and proper law.} Universal and common law only. No particular law of any diocese, eparchy, or Church \latin{sui iuris}, no proper law of any institute, and no ordinariate statute was examined. Where the universal law refers a matter outward---CIC c.~277 §3, CCEO cc.~374 and 758 §3---this article stops at the referral, and a universal capacity is never described as a current local practice.

\textbf{Material facts assumed.} None. Every norm is analyzed in the abstract; no person's status, obligations, penalty, marriage, office, or petition is assessed.

\textbf{Amendments checked.} The canons quoted were checked at the current deliveries and, for CCEO, at the promulgation text, on the as-of date. Known post-promulgation amendments material to this subject were checked: the 2021 integral revision of Book VI, which supplies the text of cc.~1394--1398 used here; and Benedict XVI's motu proprio \emph{Omnium in mentem} of 2009, whose amendments to Book IV do not touch any canon quoted here. The published collection of authentic interpretations was checked and contains none for the canons at issue. This is a check of the canons used, not a completed survey of every amending act.

\textbf{As-of date: 25 July 2026.}

\textbf{Referral.} This is a study aid, not canonical advice. Rights, penalties, marriage status, sacramental access, office, and every concrete application belong to the competent ecclesiastical authority or to a qualified canonist with the complete facts, the applicable particular law, and the current date. Section~14 sets out where particular questions go.

\subsection{Notes, rights, and review}

The numbered Notes carry exact citations, verification dates, evidentiary ceilings, and claim-local qualifications; no indispensable premise lives only in a note. Latin is quoted where the original governs; English is either quoted from an identified official delivery or supplied as a labelled working gloss. Official Holy See texts and the Internet Archive's page images of the 1918 Vatican Polyglot printing retain their own status: the former are quoted briefly with attribution from registered dated states and remain outside the project's CC~BY~4.0 grant; the latter reproduce a public-domain 1918 printing. The \emph{Acta} volume PDFs are identified by hash and retrieval route and are not redistributed, the Holy See's portal terms not granting repository redistribution. Project-created prose, tables, and organization are project content.

This revision received internal argumentative, source-consistency, quotation, rights, and production review by the authoring agent. Independent review---canonical, historical, theological, and Eastern-canonical---is outstanding; no imprimatur, nihil obstat, or ecclesiastical approval is claimed, and internal checking is not independent review. The publication language for this work is: source-audited working article.
